% RECALL -- os/samir/cmd/mutate.c m_recall (default the current record; FOR
% without a scope = ALL).  Gate: test-interp-replace.
\begin{XbaseEntry}{RECALL}{command}
Removes the deletion mark from records.
\Syntax
\SyntaxLine{\kw{RECALL} \opt{\xp{scope}} \opt{\kw{FOR} \xp{condition}}
  \opt{\kw{WHILE} \xp{condition}}}
\Defaults
Without a scope or a condition, \cmd{RECALL} works on the current record
only.  With a \kw{FOR} condition and no scope, the scope is \kw{ALL}.
\Usage
\cmd{RECALL} makes records marked by \cmd{DELETE} ordinary records again.  A
record removed by \cmd{PACK} or \cmd{ZAP} cannot be recalled.
\Examples
Record 13 of \file{STAFF.DBF} is marked for deletion on the diskette:
\DotExample{examples/c-recall}
\SeeAlso{DELETE, PACK, DELETED()}
\end{XbaseEntry}

